% Linear radix scaling, paired zero tests, and two-digit reset padding.
\section{Linear radix scaling and a certificate of 96 operations}
\label{sec:linear-radix}

The quadratic radix scale can be replaced by a linear one when the
coefficient tests are allowed two base-$B$ digits. An explicit bound
on the indicator code removes its quotient alias. Positive padding
then resets the incoming carry at every ordinary target, and a
second unit test excludes a scale of order $\sqrt B$.

Starting from Section~\ref{sec:composed-97}, replace the radix and
third packed block by
\begin{equation}\label{eq:linear-radix-block}
 B=Hb,\qquad \mathcal C=x+g,\qquad
 \Omega=\lambda-e>0,\qquad
 \sigma=\Omega(q-\mathcal C^2)>0,
\end{equation}
and add the positive gap equation
\begin{equation}\label{eq:linear-radix-gap}
 \ell+\alpha_2=q.
\end{equation}
Every Pell equation is retained, with $B=Hb$ wherever the radix
occurs. The coefficient index is rebuilt below.

\begin{theorem}\label{thm:best96}
Every recursively enumerable subset of the positive integers has
an admissible fixed index $(V,H,T_L)$ for a system with 35 positive
unknowns and 23 equations admitting a straight-line certificate of
96 operations: 52 multiplications and 44 additions. Here
$T_L=\psi_4(L)$ for the underlying exponent in the index
construction. The supplied parameters are $(x,V,H,T_L)$, with
three fixed index parameters besides the positive input $x$.
Fixed numerals and equality tests are free.
\end{theorem}

The following subsections prove the theorem. The complete arithmetic
table is in Appendix~\ref{app:96}. The exact source and primitive
checker is located in \file{../verification/}:
\begin{center}
\small\file{round30\_1980\_linear\_radix\_certificate.py}.
\end{center}

\subsection{The homogeneous circuit and its normalization rows}

Compile the represented finite Diophantine system into an
arithmetic circuit over nonnegative integers, using the actual
positive input $x$, addition, multiplication, zero and one.
Introduce a homogeneous coordinate $\delta$. Repeated operands
are replaced by fresh copies, and every gate output is fresh.
All circuit coordinates other than $x$ have positive weights in
the code.

For a multiplication, addition or zero equation use, respectively,
\begin{equation}\label{eq:linear-radix-gates}
 2XY-2W\delta=0,\qquad
 2(X+Y-W)\delta=0,\qquad
 2X\delta=0.
\end{equation}
All products in these expressions have distinct factors. A copy
$X'=X$ can instead be imposed by
\begin{equation}\label{eq:linear-radix-copy}
 X'^2-X^2=0.
\end{equation}
Nonnegativity makes this an exact copy equation. It also supplies
arbitrarily many distinct copies of $\delta$ without introducing
a diagonal coefficient of magnitude two. A circuit occurrence of
the constant one is represented by $\delta$ or a copy of $\delta$.
Whenever using $\delta$ directly would create $\delta\delta$ in
\eqref{eq:linear-radix-gates}, use a fresh square-difference copy
instead. Likewise give repeated gate operands distinct copies.
Thus every product in \eqref{eq:linear-radix-gates} has distinct
factors, including when a gate operand is the constant one. Final
equalities can use either distinct-coordinate products with
$\delta$ or square differences. Include both $F$ and $-F$ for
every zero equation $F=0$ in this construction.

For normalization, introduce distinct coordinates
$X,Y,\delta',u$ and include the following equations and their
negatives:
\begin{equation}\label{eq:linear-radix-normalization}
 \begin{gathered}
 X^2-x^2=0,\qquad Y^2-x^2=0,\qquad
 \delta'^2-\delta^2=0,\\
 2xX-2\delta\delta'-2u\delta=0.
 \end{gathered}
\end{equation}
Once their zero tests have been decoded, these equations give
\begin{equation}\label{eq:linear-radix-guard}
 X=Y=x,\qquad \delta'=\delta,\qquad
 x^2=\delta^2+u\delta.
\end{equation}
Thus $\delta>0$ because $x>0$, and $\delta\le x$.
Conversely, $\delta=\delta'=1$, $X=Y=x$ and $u=x^2-1$ are
nonnegative witnesses for every positive $x$, including $x=1$.

Finally include two unpaired targets with raw polynomial
$\delta^2$. The second will receive extra positive padding.
These two tests force $\delta=1$ below; the argument does not
presume an incoming carry at an unprotected special target.

Divide every square coefficient of a target polynomial by one
and every distinct-coordinate product coefficient by two. All
results lie in $\{-1,0,1\}$. This includes
\eqref{eq:linear-radix-gates}, \eqref{eq:linear-radix-copy},
both signs of \eqref{eq:linear-radix-normalization}, and the
two unit targets. The number of rows and encoded coordinates
may depend on the represented system; they are digits of the
integer codes, not additional unknowns of the final system.

\subsection{Two-digit support and reset positions}

Let $m$ count all non-input circuit and normalization coordinates.
Give them weights
\[
 v_i=4\,3^i\quad(0\le i<m),\qquad M=4\,3^{m-1}.
\]
The input $x$ has weight zero. Quadratic monomial weights determine
their unordered coordinate pair, including squares and products
with $x$: divide by four and read the base-three digits. All these
weights are multiples of four and are at most $2M$.

Let $s$ be the number of targets, counting both signs of each
zero equation and the two final unit tests. Put the unit tests
last, and distinguish the last target $t_*$ as the additional-carry
test. Define
\begin{equation}\label{eq:linear-radix-support}
 \begin{gathered}
 d_0=4M+4,\qquad
 t_j=(s+1)d_0+2M+jd_0\quad(0\le j<s),\\
 t_*=t_{s-1}=2sd_0+2M,\qquad K=t_*+2.
 \end{gathered}
\end{equation}
All $t_j$ are multiples of four. For a target $G_j$ and a monomial
of weight $w$ with divided coefficient $a$, put $aT^{t_j-w}$ in
$\mathcal D_{\rm main}(T)$. The intervals
$[t_j-2M,t_j]$ are disjoint, so these instructions do not collide.
Add the reset polynomial and the last-target padding:
\begin{equation}\label{eq:linear-radix-padding}
 \begin{aligned}
 \mathcal D_{\rm reset}(T)&=\sum_jT^{t_j-2},\\
 \mathcal D_{\rm extra}(T)&=T^{t_*-1}
                  +T^{t_*-1-v_X}+T^{t_*-1-v_Y},\\
 \mathcal D(T)&=\mathcal D_{\rm main}(T)
                 +\mathcal D_{\rm reset}(T)
                 +\mathcal D_{\rm extra}(T).
 \end{aligned}
\end{equation}
The three parts have exponents congruent to $0$, $2$ and $3$
modulo four, respectively. The extra exponents are distinct.
All $\mathcal D$ coefficients therefore still lie in
$\{-1,0,1\}$. Every exponent is positive, greater than $M$,
and less than $K$.

Include a dummy coordinate at each $t_j$ and $t_j+1$, and set
\begin{equation}\label{eq:linear-radix-codes}
 \begin{aligned}
 \ell_0(T)&=\sum_iT^{v_i}+\sum_j(T^{t_j}+T^{t_j+1}),\\
 e_0(T)&=\sum_{h=0}^{K-1}T^h+\mathcal D(T).
 \end{aligned}
\end{equation}
The digits of $\ell_0$ are zero or one, and those of $e_0$ are
zero, one or two. The unit digit of $e_0$ is one. Let $c_*$ count
the input, all $m$ true coordinates and the $2s$ dummies, and put
\[
 D_1=\sum_h\bigl|[T^h]\mathcal D(T)\bigr|.
\]
Choose a power of two $L>3K+2$ and a power of two $H$ such that
\begin{equation}\label{eq:linear-radix-index}
 \begin{gathered}
 H>\max\{2\,4^{2L+1},\ 64\max(1,D_1)c_*^2,\ 3L,\ 64\},\\
 V=\ell_0(4)+e_0(4)4^L,\qquad T_L=\psi_4(L).
 \end{gathered}
\end{equation}
These choices depend only on the represented system and precede
the queried input $x$.

No arbitrary decoded dummy value contributes to
$\mathcal D(T)\mathcal C(T)^2$ through position $t_*+1$.
Indeed, the least $\mathcal D$ exponent is at least $t_0-2M$,
the least dummy weight is $t_0$, and
\begin{equation}\label{eq:linear-radix-dummies}
 2t_0-2M>t_*+1.
\end{equation}
This matters because the dummies at $t_j+1$ are not in the zero
residue class modulo four. Up to the last tested digit, however,
they have no effect on the product.

\subsection{Preliminary bounds and the Pell conclusions}

Retain the coding equations and packing
\begin{equation}\label{eq:linear-radix-packing}
 \begin{aligned}
 \mathcal C&=x+g,& b&=x+\beta,& B&=Hb,&\theta&=B-4,\\
 \lambda(B-1)&=q^2-1,& n&=q^8,\\
 S_2&=\ell+eq=V+t\theta,&S_2+\alpha&=q^2,\\
 S&=g+q^2(S_2+q^2\sigma),\\
 T+1&=q^2(1+\theta\lambda)-(b-1)\ell+\theta\ell q^4,\\
 r&=S(n^2-n)+(T+1)(n^2-1).
 \end{aligned}
\end{equation}
All supplied unknowns, including $\alpha$, $\alpha_2$,
$\Omega$ and $\sigma$, are positive integers. Before any circuit
row or Pell conclusion is used, positivity gives
\begin{equation}\label{eq:linear-radix-preliminary}
 \ell<q,\quad e<q,\quad S_2<q^2,\quad
 \mathcal C^2<q,\quad g<q,\quad
 0<\sigma<\lambda q<q^3.
\end{equation}
The last inequality follows from
$\lambda=(q^2-1)/(B-1)<q^2$. Geometry gives $B\le q^2$.
Also $b\ge2$ and \eqref{eq:linear-radix-index} imply
$q\ge6$, $b<B<n$ and $3L\le B\le n$. The former estimate
$q>4b$ is not used or asserted for the linear radix. Moreover,
\[
 \theta\lambda\ge B-4=Hb-4>b.
\]
The integer bounds in \eqref{eq:linear-radix-preliminary} give
$0<S<q^7<n$. For the other packed number,
\begin{equation}\label{eq:linear-radix-packed-range}
 \begin{aligned}
 T+1&>q^2\theta\lambda-b\ell>bq^2-bq>0,\\
 T+1&<q^4(1+\theta\ell)<Bq^5\le q^7<n.
 \end{aligned}
\end{equation}
The upper bound uses $1+\theta\lambda<q^2$ and
$1+\theta\ell\le1+(B-4)(q-1)<Bq$. Consequently
\begin{equation}\label{eq:linear-radix-r-range}
 n\le n^2-1\le r<2n^3.
\end{equation}

These are all the packing-size inputs to the main Pell and
binomial arguments. Write
\[
 \begin{gathered}
 U=wn^2,\qquad Y_{\rm p}=s_{\rm p}n^2,
 \qquad a=Y_{\rm p}(U+1),\qquad A=a+4,\\
 D_{\rm p}=A^2-1,\qquad P=2UY_{\rm p}^2+1,
 \qquad J=2r+1.
 \end{gathered}
\]
The notation $Y_{\rm p},s_{\rm p}$ distinguishes the Pell values
from the encoded input copy $Y$. We have
$U,Y_{\rm p}\ge n^2$, $UY_{\rm p}\ge n^4>r+1$ and
$a>n^4>J$. The first index equation and positive interval give
$c>J$. The norm classifications and $P>A$ imply that the preliminary
main index is at least $r+2$, hence $c>AD_{\rm p}^2$, exactly as
in Section~\ref{sec:composed-97}.

The retained relaxed auxiliary norm
\[
 (ic^2)^2=D_{\rm p}(f^2-1)
\]
therefore has the same fundamental-unit and strong-divisibility
hypotheses. It recovers the required auxiliary multiple of $c$.
The doubled signed-index argument identifies
$c=\psi_A(J)$ and $d=\chi_A(J)$. The first-index and ratio
argument gives $Y_{\rm p}\ge U^r$ and $a>U^{r+1}$ before either
exponent relation is used. The common-witness exponent equation
then gives $U=4^J$. The retained fixed congruence
$\kappa=T_L+\Delta a$ determines the second index as $L$, and
the second exponent equation gives $q=B^L$.

Every exponent-size hypothesis follows from
\eqref{eq:linear-radix-r-range} and $3L\le B\le n\le r$ as
in the preceding proof. None requires $B=Hb^2$. Since $n^2$ divides
$U=4^J$, $q$ is a power of two; $q=B^L$ makes $B$ a power of two;
and the fixed power of two $H$ divides $B=Hb$, so $b$ is a power
of two. The unchanged upper ratio and binomial fractional-part
estimates give
\begin{equation}\label{eq:linear-radix-binomial}
 n^2\mid\binom{2r}{r}.
\end{equation}
This applies the full Pell argument before any coefficient
equation or unit value is assumed. Its necessity construction
will be used only after the new coding witnesses have been chosen.

\subsection{Canonical codes without a high-mask estimate}

After the Pell conclusions, $b<B<q$ and $\ell<q$. Therefore
\begin{equation}\label{eq:linear-radix-no-borrow}
 q^2-1-(b-1)\ell>0.
\end{equation}
The packed number $T$ now has the exact three-block expansion
\[
 T=\bigl(q^2-1-(b-1)\ell\bigr)
                +q^2\theta\lambda+q^4\theta\ell.
\]
The blocks have widths $(2,2,4)$ in radix $q$. The binary
no-carry packing lemma, using
\eqref{eq:linear-radix-packing}, \eqref{eq:linear-radix-r-range}
and \eqref{eq:linear-radix-binomial}, gives
\begin{equation}\label{eq:linear-radix-masks}
 \begin{gathered}
 g\mathbin{\&}\bigl(q^2-1-(b-1)\ell\bigr)=0,\\
 S_2\mathbin{\&}\theta\lambda=0,\qquad
 \sigma\mathbin{\&}\theta\ell=0.
 \end{gathered}
\end{equation}
Here $\mathbin{\&}$ denotes bitwise intersection only in the
proof of the polynomial system's meaning.

Since $q=B^L$ and $\theta=B-4$, the middle mask restricts the
base-$B$ digits of $S_2$ to $\{0,1,2,3\}$ in fewer than $2L$
positions. If $F(T)$ is its digit polynomial, then $F(B)=S_2$,
and the congruence in \eqref{eq:linear-radix-packing} gives
$F(4)\equiv V\pmod{B-4}$. Both $F(4)$ and $V$ are less than
$4^{2L}$, while \eqref{eq:linear-radix-index} makes $B-4$
larger than their possible difference. Hence $F(4)=V$.
Uniqueness of the base-four expansion gives
\[
 S_2=\ell_0(B)+e_0(B)q.
\]
Both $\ell,e$ and $\ell_0(B),e_0(B)$ are less than $q$.
Uniqueness of quotient and remainder on division by $q$ yields
\begin{equation}\label{eq:linear-radix-canonical}
 \ell=\ell_0(B),\qquad e=e_0(B).
\end{equation}
No high-mask estimate proportional to the square of the digit
bound is required.

The first mask, together with $b$ and $B$ being powers of two,
now bounds every coordinate digit of $g$ by $b-1$ and permits
nonzero digits only at the indicator positions. Its unit digit
is zero, so the unit coordinate of $\mathcal C$ is the queried
$x$. Write
\[
 \mathcal C(T)=x+\sum_i z_iT^{v_i}
              +\sum_j(d_jT^{t_j}+d'_jT^{t_j+1}).
\]
All decoded coordinates are nonnegative and below $b$. Put
$A_{\mathcal C}=\mathcal C(1)^2$. Then
\begin{equation}\label{eq:linear-radix-coefficient-bound}
 A_{\mathcal C}<(c_*b)^2<B^2/64.
\end{equation}
Also $e,\mathcal C<B^K$, and $L>3K+2$ gives
$e\mathcal C^2<q$. Through position $K-1$, the raw signed
coefficients of $\sigma$ are those of
$\mathcal D(T)\mathcal C(T)^2$: the $q$ terms in
\eqref{eq:linear-radix-block} start at $L$, and the difference
between the finite baseline and $\lambda$ starts at $K$.
Each raw coefficient has absolute value at most $A_{\mathcal C}$,
because the coefficients of $\mathcal D$ have absolute value
at most one.

\subsection{Reset padding and exact paired zero tests}

Consider the signed base-$B$ carry recursion
\begin{equation}\label{eq:linear-radix-carry}
 \begin{gathered}
 h_0=0,\qquad a_j=[T^j]\mathcal D(T)\mathcal C(T)^2,\\
 h_{j+1}=\left\lfloor\frac{a_j+h_j}{B}\right\rfloor,
 \qquad \mathrm{digit}_j=a_j+h_j-Bh_{j+1}.
 \end{gathered}
\end{equation}
It computes the actual low digits of the positive integer
$\sigma$. Bound \eqref{eq:linear-radix-coefficient-bound}
gives $|h_j|<B/8$ by induction. Indeed,
\[
 \left|\left\lfloor\frac{a_j+h_j}{B}\right\rfloor\right|
 \le \frac B{64}+\frac18+1<\frac B8
\]
when $|a_j|<B^2/64$, $|h_j|<B/8$ and $B>64$.

By \eqref{eq:linear-radix-dummies}, dummies do not affect the
tested range. All true non-input weights are zero modulo four.
The raw coefficients can therefore be nonzero only in residue
classes $0$, $2$ and $3$ modulo four; class $1$ is empty. The
carry into $t_j-2$, following an empty class-$1$ position, is
consequently either $-1$ or zero.

At $t_j-2$, the raw coefficient is nonnegative and at least
$x^2$: the row's own reset contributes $x^2$, and every other
reset contribution is nonnegative. It is less than $B^2/64$.
The carry into $t_j-1$ is therefore a nonnegative integer
$\alpha_j<B/64$.

At every target except $t_*$, the raw coefficient at $t_j-1$
is zero. Only $\mathcal D_{\rm extra}$ contributes in this
residue class, and its contributions are within distance $2M$
of $t_*-1$; different targets are separated by $4M+4$.
Dividing $\alpha_j<B$ at this empty position leaves exactly
zero incoming carry at $t_j$.

There is no raw coefficient at $t_j+1$. Thus the two tested
digits at $t_j,t_j+1$ form the residue modulo $B^2$ of the raw
target polynomial $G_j$. The third mask permits both digits
to lie only in $\{0,1,2,3\}$.

For a paired zero equation, $|G_j|<B^2/64$. A negative nonzero
value has a residue modulo $B^2$ whose high digit exceeds three,
so it fails the mask. Both $G_j$ and $-G_j$ have their own
reset and zero incoming carry. They can both pass exactly when
$G_j=0$. This proves every circuit and normalization equation
exactly, before $\delta$ is known to be one. In particular
\eqref{eq:linear-radix-guard} holds.

The first unpaired $\delta^2$ target also receives zero incoming
carry. The sole exception to this rule is the deliberately
modified last target.

\subsection{The last carry fixes the unit}

The first unit test has raw value $\delta^2$ and zero incoming
carry. Its two allowed digits give
\begin{equation}\label{eq:linear-radix-unit-digits}
 \delta^2=kB+s,\qquad k,s\in\{0,1,2,3\}.
\end{equation}
The guard gives $0<\delta\le x$, and the coordinate bound gives
$\delta<b=B/H<B/4$. Since $B$ is a sufficiently large power
of two, $s=2$ or $3$ is impossible modulo four. If $s=1$, the
square roots of one modulo $B$ are congruent to
$1,B/2-1,B/2+1,B-1$. The bound $\delta<B/4$ forces
$\delta=1$. The only remaining case for $\delta>1$ is $s=0$,
and therefore
\begin{equation}\label{eq:linear-radix-large-unit}
 \delta^2\ge B,\qquad x^2\ge B.
\end{equation}

At $t_*-1$, monomial uniqueness and the input-copy equations
make the raw extra-padding coefficient exactly
\[
 x^2+2xX+2xY=5x^2.
\]
The reset at $t_*-2$ remains present and supplies a nonnegative
carry $\alpha_*<B/64$. Thus the incoming carry at the last
target is
\begin{equation}\label{eq:linear-radix-unit-carry}
 k_{\rm pad}=\left\lfloor\frac{5x^2+\alpha_*}{B}\right\rfloor
 \ge0.
\end{equation}
The global coefficient bound also gives $k_{\rm pad}<B/8<B$.
If \eqref{eq:linear-radix-large-unit} held, then
$k_{\rm pad}\ge5$. Since $\delta^2$ is a multiple of $B$,
the low digit at the last target would be $k_{\rm pad}$,
between five and $B-1$, contradicting its mask. Therefore
$\delta=1$.

The circuit equations now recover the represented system at the
actual positive input $x$. This proves sufficiency.

\subsection{Necessity and all positive supplied witnesses}

Given a solution of the represented system, take its circuit
values, $\delta=\delta'=1$, $X=Y=x$, $u=x^2-1$, all fresh
copies equal, and every dummy zero. Each paired row has value
zero, and both unpaired unit targets have value one.

The fixed index was already chosen in
\eqref{eq:linear-radix-support}--\eqref{eq:linear-radix-index}.
Now choose a power of two $b$ arbitrarily large relative to the
finitely many coordinate values. Besides exceeding every
coordinate, choose it so that every raw coefficient of
$\mathcal D(T)\mathcal C(T)^2$ has absolute value less than
$B/4$ for $B=Hb$, and $5x^2<B/4$. This is possible because
the coordinates, $\mathcal D$ and their raw coefficients are
fixed before $b$ is chosen, whereas $Hb$ grows without bound.

Take the canonical codes, $q=B^L$ and $n=q^8$. The support
margin gives $\mathcal C^2<q$, and
$\lambda=(q^2-1)/(B-1)>e$. Hence $\Omega$ and $\sigma$ are
positive. The values
\[
 \begin{gathered}
 \alpha_2=q-\ell,\qquad
 \alpha=q^2-(\ell+eq),\qquad \beta=b-x,\\
 t=\frac{\ell+eq-V}{B-4}
 \end{gathered}
\]
are positive integers. Positivity of $t$ follows by evaluating
the nonconstant nonnegative polynomial
$\ell_0(T)+e_0(T)T^L$ at $B>4$. The coordinate $\delta=1$
ensures $g>0$.

The first two masks hold by the coordinate and digit bounds.
For the third, all variable positions and their preceding
positions lie below the least $\mathcal D$ exponent, so their
raw coefficients and carries are zero. Every paired target has
zero value and receives zero incoming carry from its reset.
The first unit target has value one. At the last target, the
chosen small-coefficient bounds give $\alpha_*=0$ and
$k_{\rm pad}=0$, so its tested digits are also one and zero.
Thus all masks hold. The packing lemma gives
\eqref{eq:linear-radix-binomial} for the newly determined $r$,
with the required bounds already established.

Construct the positive Pell witnesses at this new $r$ as in
Section~\ref{sec:composed-97}: choose $U=4^J$,
\[
 Y_{\rm p}=\left\lfloor\frac{(U+1)^{2r}}{U^r}\right\rfloor,
 \qquad w=\frac U{n^2},\qquad
 s_{\rm p}=\frac{Y_{\rm p}}{n^2},\qquad a=Y_{\rm p}(U+1).
\]
Take the main and first Pell coordinates at $J$ and $r+1$,
their positive interval and first-index quotients, and the
second coordinates at $L$. The second index uses the newly
fixed $T_L=\psi_4(L)$. The exponent and rounding bounds give
every required positive quotient as in that proof.

Choose the doubled-index auxiliary exponent last as
$m_{\rm aux}=2cJ$. Set $f=\chi_A(m_{\rm aux})$ and
$i_{\rm old}=\psi_A(m_{\rm aux})/c^2>0$; the doubled-index
construction supplies positive $o,j$. Use
$i=D_{\rm p}i_{\rm old}$. This gives the relaxed norm
$(ic^2)^2=D_{\rm p}(f^2-1)$ and preserves the signed auxiliary
value. No coding witness changes. Together with the new positive
$\alpha_2$, all 35 supplied unknowns are positive integers.
This proves necessity and completes the mathematical proof of
Theorem~\ref{thm:best96}.

\subsection{The complete arithmetic count}

Relative to the 97-operation certificate, replacing $B=Hb^2$
by $Hb$ removes one multiplication. Replacing
\[
 \sigma=\theta\lambda q-(\lambda-e)\mathcal C^2
 \quad\hbox{by}\quad
 \sigma=(\lambda-e)(q-\mathcal C^2)
\]
replaces two multiplications and one subtraction by one
multiplication and one subtraction, removing another
multiplication. Equation~\eqref{eq:linear-radix-gap} costs
one addition. All normalization, padding and two-digit support
changes belong to the fixed index rather than to the arithmetic
of the final system. Therefore
\[
 97-1-1+1=96,\qquad
 (54\text{ multiplications},43\text{ additions})
 \longrightarrow(52\text{ multiplications},44\text{ additions}).
\]

The standalone checker verifies all 96 primitive instructions,
the complete 23 source residuals and their acyclic corrections.
The geometric and second-exponent corrections depend only on
the exact equality $\theta+4=B$. The signed auxiliary
correction depends only on the exact relaxed auxiliary norm.
The sparse coefficient/carry regression supplies additional
finite evidence; it does not replace the support, carry, index
or positive-witness arguments above. The result is an upper
bound, not an optimality claim.
